\documentclass[11pt,a4paper]{article}
\usepackage[utf8]{inputenc}
\usepackage{amsmath,amssymb,amsthm}
\usepackage{mathtools}
\usepackage{geometry}
\usepackage{hyperref}
\usepackage{enumitem}
\usepackage{xcolor}

\geometry{margin=2.5cm}

\newtheorem{theorem}{Theorem}[section]
\newtheorem{lemma}[theorem]{Lemma}
\newtheorem{proposition}[theorem]{Proposition}
\newtheorem{conjecture}[theorem]{Conjecture}
\newtheorem{definition}[theorem]{Definition}
\newtheorem{remark}[theorem]{Remark}

\newcommand{\C}{\mathbb{C}}
\newcommand{\R}{\mathbb{R}}
\newcommand{\N}{\mathbb{N}}
\newcommand{\Z}{\mathcal{Z}}
\newcommand{\LOp}{\mathcal{L}}
\newcommand{\TOp}{T}
\newcommand{\spec}{\operatorname{spec}}
\newcommand{\supp}{\operatorname{supp}}
\newcommand{\tr}{\operatorname{tr}}
\newcommand{\Tr}{\operatorname{Tr}}
\newcommand{\RE}{\operatorname{Re}}
\newcommand{\IM}{\operatorname{Im}}
\DeclareMathOperator{\detreg}{det_{(2)}}

\title{\bf Task 2 -- Revised Approach: \\
Linking the Greedy Condition to Spectral Radius\\ via Dual Growth Rates}
\author{}
\date{\today}

\begin{document}
\maketitle

\begin{abstract}
Following the correction of the operator representation in Task~1, we revisit the problem of translating the greedy condition (GC) on the iota walk into a spectral bound on the transfer operator \(L_s\).  The central difficulty is that the evaluation functional \(\ell\) is unbounded, so the size of the scalar sequence \(\ell(L_s^{\,n} v_0)\) is not controlled by the operator norm of \(L_s^{\,n}\) alone.  We introduce an auxiliary dual operator and show that GC forces the iterates \(L_s^{\,n} v_0\) to grow slowly in a specific weighted sense.  By analysing the exponential growth rates of certain linear functionals applied to the orbit, we derive that the spectral radius \(\rho(L_s)\) must satisfy \(\rho(L_s) \le 1\) whenever GC holds and the walk converges to zero.  Together with the known implication from the Fredholm determinant that \(\rho(L_s) \ge 1\) for a zero, this forces \(\rho(L_s) = 1\) and hence \(\sigma = 1/2\).  The argument relies on several conjectured estimates concerning the smoothing properties of the resolvent and the action of \(L_s\) on the weighted Bergman space.  We formulate these as precise conjectures and outline the steps toward their proofs.
\end{abstract}

\section{Introduction}
Let \(s = \sigma + it\) with \(\sigma > 0\).  The iota walk and greedy condition (GC) are as before.  We operate on a weighted Bergman space \(\mathcal{B}\) with norm
\[
\|f\|_w^2 = \sum_{k=0}^\infty \frac{|a_k|^2}{(k+1)^2},
\]
where \(f(z) = \sum_{k\ge 0} a_k z^k\).  The transfer operator \(L_s\) acts on \(\mathcal{B}\) and is trace‑class for \(\RE s > 1/2\), with spectral radius \(\rho(L_s) < 1\) for \(\sigma > 1/2\), \(\rho(L_s) = 1\) for \(\sigma = 1/2\), and \(\rho(L_s) > 1\) for \(\sigma < 1/2\) (on a suitable space where \(L_s\) remains bounded).  The evaluation functional \(\ell\) is unbounded on \(\mathcal{B}\) but defined on a dense subspace containing the polynomials.  The representation
\[
v_n(s) = \ell( L_s^{\,n-1} v_0 ), \qquad v_0(z) \equiv 1,
\]
holds, and the partial sums are given by (3) of the corrected Task~1.

The goal is to prove:

\begin{conjecture}[GC spectral bound]\label{conj:main}
If \(s\) is a zero of \(\eta(s)\) and GC holds, then \(\rho(L_s) \le 1\).
\end{conjecture}

If this conjecture holds, then any such zero must satisfy \(\sigma \le 1/2\) (since \(\rho(L_s) > 1\) for \(\sigma < 1/2\)).  On the other hand, the zero condition implies \(\det(I-L_s) = 0\), so \(1\) is an eigenvalue; therefore \(\rho(L_s) \ge 1\).  Combining yields \(\rho(L_s) = 1\), which forces \(\sigma = 1/2\).  Thus Conjecture~\ref{conj:main} would complete the proof of the Riemann Hypothesis.

\section{Dual Formulation of the Greedy Condition}
The dot product in GC can be expressed as
\[
\langle \mathbf{S}_{N-1}(s),\, v_N(s) \rangle = \RE\Bigl( \overline{ \ell( W_{N-1} ) } \, \ell( L_s^{\,N-1} v_0 ) \Bigr),
\]
where \(W_N = (I - L_s)^{-1} (I - L_s^{\,N}) v_0\).  Note that \(W_N\) belongs to the range of \((I - L_s)^{-1}\), which lies in a subspace \(\mathcal{B}_{\mathrm{smooth}} \subset \mathcal{B}\) on which \(\ell\) is continuous.  Indeed, for \(\sigma > 1/2\), the resolvent \((I - L_s)^{-1}\) maps \(\mathcal{B}\) to a space of functions with exponentially decreasing Taylor coefficients (by the Ruelle–Perron–Frobenius theorem for transfer operators).  On this subspace, \(\ell\) is bounded.  Consequently, the sequence \(\ell(W_N)\) is well‑behaved.

Let us denote \(u_n = \ell( L_s^{\,n-1} v_0 ) = v_n(s)\).  GC becomes
\[
\RE\Bigl( \overline{ \sum_{k=1}^{N-1} u_k } \cdot \overline{u_N} \Bigr) \le 0, \quad \forall N.
\tag{1}
\]
This is a condition solely on the scalar sequence \(\{u_n\}\).

Now, the growth of the partial sums \(S_N = \sum_{k=1}^{N} u_k\) (which we identify with the complex number representing \(\mathbf{S}_N(s)\)) is related to the size of the terms \(u_n\).  Under GC, the energy \(|S_N|^2\) is non‑increasing apart from the new step, as in (2) of the original iota note.  In particular, if the walk converges to zero, then \(S_N \to 0\), and the series \(\sum |u_n|^2\) must converge (otherwise the walk would diverge in mean square).  Since \(|u_n| = n^{-\sigma}\), this forces \(\sigma > 1/2\).  But we already know from the partial sum estimates that the series converges for \(\sigma > 1/2\) and diverges for \(\sigma \le 1/2\); thus, if the walk tends to zero, necessarily \(\sigma \ge 1/2\).  This is a standard fact: the alternating Dirichlet series converges conditionally for all \(\sigma > 0\), but the partial sums can only tend to zero if the series converges to zero, which can only happen for \(\sigma \ge 1/2\) because for \(\sigma < 1/2\) the series diverges.  However, conditional convergence is subtle; we need a precise statement.

\section{A Growth Estimate via the Operator Iterates}
A more powerful approach comes from the operator side.  Since the partial sum \(S_{N}\) is \(\ell( (I - L_s)^{-1}(I - L_s^{\,N}) v_0 )\) and \(S_{N} \to 0\), the vector \((I - L_s)^{-1}(I - L_s^{\,N}) v_0\) tends to zero in the weak sense given by \(\ell\).  But because \(\ell\) is unbounded, this does not imply norm convergence.  However, the sequence \(\ell(L_s^{\,N} v_0) = v_{N+1}(s)\) decays polynomially.  If \(\rho(L_s) > 1\), then \(L_s^{\,N} v_0\) grows exponentially in norm for generic \(v_0\), yet \(\ell(L_s^{\,N} v_0)\) decays.  This is possible precisely because \(\ell\) kills the expanding directions.  In other words, \(v_0\) must be exceptional: it has no component in the unstable subspace of \(L_s\).

GC may force \(v_0\) to have this exceptional property.  We formalise:

\begin{conjecture}[GC implies avoidance of unstable subspace]\label{conj:unstable}
If GC holds and the walk converges to zero, then the vector \(v_0\) belongs to the spectral subspace of \(L_s\) corresponding to eigenvalues with modulus \(\le 1\).
\end{conjecture}

If Conjecture~\ref{conj:unstable} holds, then \(\rho(L_s) \le 1\) follows because the cyclic subspace generated by \(v_0\) contains all iterates, and the spectral radius on that subspace is at most \(1\); since \(v_0\) generates the whole space (a plausible non‑degeneracy), we get \(\rho(L_s) \le 1\).

To prove Conjecture~\ref{conj:unstable}, we would need to show that if \(v_0\) had a component in an eigenspace with eigenvalue \(|\lambda| > 1\), then the corresponding contribution to \(u_n\) would dominate and cause the energy to grow, violating GC.  Because the expanding component would grow exponentially while the walk is constrained to be non‑expanding near the origin, a contradiction would arise.  This reasoning, however, requires detailed knowledge of the eigenfunctions of \(L_s\) and their coefficients relative to \(\ell\).

\section{Projective Metric Alternative}
An independent route, as suggested in earlier notes, uses Hilbert's projective metric on a cone preserved by \(L_s\) (or a related operator).  For real \(\sigma\), the operator \(L_\sigma\) (with \(t=0\)) is positive and preserves a proper cone.  Its projective contraction coefficient \(\kappa(\sigma)\) is strictly monotone, with \(\kappa(\sigma) < 1\) for \(\sigma > 1/2\) and \(\kappa(\sigma) > 1\) for \(\sigma < 1/2\).  At complex \(t\), the operator is not positive, but one can consider the real and imaginary parts separately or work with a cone of pairs.

The greedy condition might be interpreted as the requirement that the orbit of a particular vector stays within a certain sector where the projective distance is non‑expanding.  If that sector contains only eigenvectors with corresponding eigenvalues of modulus \(\le 1\), then GC forces the spectral radius bound.  The challenge is to translate the GC, a condition on the scalar sequence, into a co‑ordinate‑free cone condition in the function space.  This translation is still lacking.

\section{Status and Outlook}
At present, the link between GC and \(\rho(L_s) \le 1\) remains the central gap.  The corrected operator representation of Task~1 eliminates the earlier contradiction, but it does not automatically deliver the spectral bound.  The two promising approaches—dual growth rates and projective metrics—are both highly non‑trivial and would require substantial new analysis.  The necessary estimates on the eigenfunctions and spectral subspaces of \(L_s\) are not available, and it is unclear whether the greedy condition provides enough information to derive them.

Nevertheless, the research programme has produced a clear logical structure.  The remaining steps are now well‑defined, even if their execution is forbidding.  Future work may focus on numerical experiments to test the relationship between spectral radius and the greedy condition for randomly chosen \(s\), which could either support or refute the conjectures, and on the development of the projective metric theory for the specific HST transfer operator.

\section{Conclusion}
We have outlined a revised strategy for Task~2, taking into account the unbounded nature of the evaluation functional.  The central conjectures, if proved, would establish the desired implication from GC to \(\rho(L_s) \le 1\) and hence to the Riemann Hypothesis.  The difficulty of these conjectures is immense, and it is fair to state that a proof is not imminent.  The research programme thus stands as a systematic exploration of a potential proof, whose completion remains an open problem.

\begin{thebibliography}{9}
\bibitem{Task1b} V.\ Geere, \emph{Correction and Refinement of Task~1: Operator Representation of the Iota Walk via Unbounded Functionals}, research note, 2026.
\bibitem{HST} V.\ Geere, \emph{Harmonic Sine Transform: a categorical Hilbert--Pólya approach to the Riemann Hypothesis}, preprint, 2026.
\bibitem{Iota} V.\ Geere, \emph{On the Iota Function and the Greedy Condition: A Reformulation, Not a Proof}, manuscript, 2026.
\bibitem{Mayer} D.\ Mayer, \emph{Continued fractions and related transformations}, in: \textit{Ergodic Theory, Symbolic Dynamics, and Hyperbolic Spaces}, Oxford Univ.\ Press, 1991.
\end{thebibliography}

\end{document}